\chapter{Résultats et évaluation}
\label{chap:resultats}

\section{Protocole expérimental}

L'évaluation de \logminer{} repose sur plusieurs catégories de données. Ce choix permet d'observer le comportement du prototype face à des sources qui diffèrent par leur structure, leur origine et leur niveau de labellisation. Les expérimentations utilisent ainsi des journaux collectés localement, des exports provenant d'un SIEM/HIDS, des jeux de données publics de journaux système, des flux réseau labellisés ainsi que plusieurs entrées spécifiquement préparées pour tester la robustesse du pipeline.

\begin{table}[H]
\centering
\caption{Jeux de données et scénarios expérimentaux}
\label{tab:datasets}
\scriptsize

\begin{tabularx}{\textwidth}{p{2.4cm}p{3.2cm}p{3cm}X}
\toprule
\textbf{Catégorie} &
\textbf{Source} &
\textbf{Nature} &
\textbf{Objectif}
\tabularnewline
\midrule

Réel local &
Windows Event/Application/System &
Journaux Windows &
Collecte, parsing et détection locale.
\tabularnewline

Réel/SIEM &
Exports Wazuh &
Alertes et événements SIEM/HIDS &
Détection non supervisée appliquée à des signaux de sécurité.
\tabularnewline

Réel tabulaire &
Linux/auth &
Données d'authentification labellisées &
Évaluation supervisée et analyse des faux positifs.
\tabularnewline

Public système &
HDFS, BGL &
Journaux distribués/HPC &
Évaluation exploratoire de journaux à forte composante séquentielle.
\tabularnewline

Public réseau &
CICIDS2017, CIC-DDoS2019 &
Flux réseau labellisés &
Détection supervisée d'attaques connues à l'aide de protocoles exploratoires et de partitions plus strictes.
\tabularnewline

Contrôle &
Windows simulé &
Événements comportant des anomalies injectées &
Comparaison entre une méthode de référence et une approche fondée sur l'apprentissage automatique.
\tabularnewline

Robustesse &
Apache, CEF/LEEF, CloudTrail, Linux auth &
Multiformat &
Évaluation de la tolérance du pipeline à plusieurs structures de journaux.
\tabularnewline

\bottomrule
\end{tabularx}
\end{table}

Cette diversité permet d'éviter une évaluation limitée à un seul type de données. Elle ne signifie cependant pas que toutes les expériences répondent à la même question. Les données labellisées servent principalement à mesurer les performances de classification, alors que les données non labellisées permettent plutôt d'étudier la production d'anomalies candidates et le fonctionnement de la chaîne d'analyse.

\section{Configuration expérimentale des modèles}

Le tableau \ref{tab:model-config} rassemble les principaux paramètres utilisés pour les modèles supervisés et non supervisés. Il fournit une vue synthétique des configurations retenues. Les scripts conservés avec le projet restent la référence détaillée pour la reproduction exacte des expériences.

\begin{table}[H]
\centering
\caption{Configuration compacte des modèles}
\label{tab:model-config}
\scriptsize

\begin{tabularx}{\textwidth}{p{2.2cm}p{2.1cm}p{1.5cm}p{1.5cm}p{1.4cm}X}
\toprule
\textbf{Famille} &
\textbf{Modèle} &
\textbf{n\_estimators} &
\textbf{max\_depth} &
\textbf{Graine} &
\textbf{Partition / paramètres clés}
\tabularnewline
\midrule

Linux/auth &
Random Forest &
120 &
24 &
42 &
Partition stratifiée exploratoire ; \texttt{class\_weight=balanced}, \texttt{min\_samples\_leaf=2}.
\tabularnewline

CICIDS2017 &
Random Forest &
100 &
28 &
42 &
Partition stratifiée exploratoire ; \texttt{class\_weight=balanced}, \texttt{min\_samples\_leaf=2}.
\tabularnewline

CIC-DDoS2019 &
Random Forest &
100 &
28 &
42--46 &
Holdout par fichier ou famille DDoS ; \texttt{class\_weight=balanced}, échantillons plafonnés.
\tabularnewline

Ablation modèle global/familles &
Random Forest &
120 &
22 &
42--46 &
Espace commun strict, partition 75/25 par famille, PR-AUC et MCC rapportés.
\tabularnewline

Validation stricte &
Random Forest &
100--120 &
24--28 &
42--46 &
Holdout par serveur, fichier ou scénario.
\tabularnewline

Windows/Wazuh &
Isolation Forest &
200 &
-- &
42 &
\texttt{contamination=0,05} ; les sorties sont interprétées comme des anomalies candidates et non comme des intrusions confirmées.
\tabularnewline

HDFS/BGL Drain3 &
Isolation Forest &
200 &
-- &
42 &
Contamination fixée sur l'entraînement ; modèles de messages Drain3 et fenêtres temporelles.
\tabularnewline

\bottomrule
\end{tabularx}
\end{table}

Les paramètres ont été choisis de manière à conserver des modèles relativement légers et reproductibles dans un environnement local. L'objectif n'est pas de rechercher systématiquement la meilleure configuration possible pour chaque jeu de données, mais de disposer d'un cadre suffisamment stable pour comparer les comportements observés selon les familles et les protocoles de séparation.

\section{Résultats principaux}

Le tableau \ref{tab:main-results} regroupe les résultats considérés comme les plus représentatifs des différentes familles étudiées.

\begin{table}[H]
\centering
\caption{Résultats principaux}
\label{tab:main-results}
\scriptsize

\begin{tabularx}{\textwidth}{p{2.2cm}p{2.8cm}p{2.3cm}p{2cm}X}
\toprule
\textbf{Famille} &
\textbf{Jeu de données} &
\textbf{Modèle} &
\textbf{Type} &
\textbf{Résultat}
\tabularnewline
\midrule

Windows &
Security/Application/System &
Isolation Forest &
Non supervisé &
Production d'anomalies candidates et d'incidents corrélés.
\tabularnewline

Wazuh &
Exports Jan./Oct./Déc. &
Isolation Forest &
Non supervisé &
122 563 événements traités et 3 676 anomalies candidates.
\tabularnewline

Linux/auth &
\texttt{linux\_auth\_logs} &
Random Forest &
Supervisé &
F1 = 0,916602.
\tabularnewline

CICIDS2017 &
MachineLearningCVE &
Random Forest &
Supervisé &
F1 exploratoire = 0,997163 ; F1 strict = 0,157827.
\tabularnewline

CIC-DDoS2019 &
DrDoS/UDPLag/SYN &
Random Forest &
Supervisé strict &
F1 = 0,997577 ; MCC = 0,994091.
\tabularnewline

\bottomrule
\end{tabularx}
\end{table}

Ces valeurs montrent immédiatement que la lecture d'un score doit rester liée au protocole utilisé. Le cas de CICIDS2017 est particulièrement parlant : un même type de modèle peut obtenir un score très élevé avec une séparation aléatoire, puis perdre une grande partie de ses performances lorsque des scénarios entiers sont exclus de l'entraînement.

\section{Évaluation des modèles supervisés}

Les modèles supervisés sont évalués sur des données pour lesquelles des labels sont disponibles. Les premiers protocoles utilisent des partitions exploratoires, généralement stratifiées. Ces expériences montrent qu'un modèle spécialisé peut très bien séparer les classes lorsqu'il est entraîné et testé sur des observations provenant de distributions proches.

Cette lecture doit toutefois rester prudente. Une partition aléatoire effectuée au niveau des lignes peut conduire des échantillons provenant d'une même capture, d'un même serveur ou d'un même scénario à se retrouver simultanément dans les ensembles d'entraînement et de test. Dans ce cas, un score élevé peut refléter en partie la similarité entre les deux ensembles plutôt qu'une véritable capacité de généralisation vers une situation nouvelle.

Le score de 0,999965 historiquement obtenu sur des données DDoS est attribué à un jeu compatible avec CIC-DDoS2019. Il ne correspond pas à une validation stricte sur UNSW-NB15. Ce résultat est conservé comme repère exploratoire, mais il n'est plus considéré comme la preuve principale du comportement du modèle sur les attaques DDoS.

Le résultat retenu pour CIC-DDoS2019 provient plutôt du protocole de holdout par fichier ou par famille DDoS, exécuté sur cinq graines. Il produit un F1 de 0,997577, un PR-AUC de 0,999991 et un MCC de 0,994091. Ces valeurs sont conservées dans le fichier \texttt{random\_forest\_cic\_ddos2019\_metrics.csv}.

Une expérimentation réellement consacrée à UNSW-NB15 devrait être menée séparément à partir des fichiers officiels de ce jeu de données. Elle devrait clairement documenter les ensembles d'entraînement et de test et produire ses propres métriques, sans réutiliser les résultats obtenus sur CIC-DDoS2019.

\begin{table}[H]
\centering
\caption{Évaluation supervisée stricte complémentaire}
\label{tab:supervised-strict-splits}
\scriptsize

\begin{tabularx}{\textwidth}{p{2.2cm}Xccccc}
\toprule
\textbf{Jeu de données} &
\textbf{Partition} &
\textbf{Rép.} &
\textbf{Précision} &
\textbf{Rappel} &
\textbf{F1} &
\textbf{MCC}
\tabularnewline
\midrule

Linux/auth &
Holdout par serveur &
5 &
0,887734 &
0,774720 &
0,826251 &
0,683796
\tabularnewline

CICIDS2017 &
Cinq holdouts par fichier ou scénario &
5 &
0,390476 &
0,128650 &
0,157827 &
0,146206
\tabularnewline

CIC-DDoS2019 &
Holdout par fichier ou famille DDoS &
5 &
0,999598 &
0,995600 &
0,997577 &
0,994091
\tabularnewline

\bottomrule
\end{tabularx}
\end{table}

Cette seconde série d'expériences complète donc les partitions exploratoires avec des séparations plus exigeantes.

Pour maintenir un coût de calcul raisonnable sur les fichiers volumineux, les échantillons sont plafonnés. Les modèles réseau utilisent 100 arbres, tandis que le modèle Linux/auth en utilise 120.

Sur Linux/auth, les performances diminuent mais restent relativement stables : le F1 passe à 0,826251. La situation est beaucoup plus marquée sur CICIDS2017, où le F1 moyen tombe à 0,157827 lorsque des fichiers ou scénarios entiers sont exclus de l'entraînement.

Le rappel n'est alors que de 0,128650. Autrement dit, une grande partie des événements positifs appartenant aux scénarios jamais vus pendant l'entraînement ne sont pas reconnus par le modèle.

CIC-DDoS2019 conserve au contraire un résultat élevé dans le protocole retenu. Cette performance reste néanmoins limitée aux familles DDoS évaluées. Elle ne permet pas d'étendre automatiquement la conclusion à d'autres catégories d'attaques ou à d'autres jeux de données.

Afin de mieux isoler l'effet de la stratégie de séparation sur CICIDS2017, une expérience contrôlée supplémentaire a été réalisée. Les caractéristiques utilisées, le modèle Random Forest, les hyperparamètres, les graines ainsi que les tailles d'entraînement et de test restent appariés. La principale différence entre les deux configurations concerne donc la manière de construire les ensembles d'entraînement et de test.

La première utilise une séparation aléatoire stratifiée. La seconde maintient chaque fois un fichier ou un scénario complet hors de l'entraînement.

\begin{table}[H]
\centering
\caption{Comparaison contrôlée entre partition aléatoire et holdout sur CICIDS2017}
\label{tab:controlled-split-cicids}
\scriptsize

\begin{tabularx}{\textwidth}{p{3.0cm}p{1.2cm}p{2.0cm}p{2.2cm}p{2.2cm}X}
\toprule
\textbf{Partition} &
\textbf{Rép./holdouts} &
\textbf{Test moy.} &
\textbf{F1 $\mu \pm \sigma$} &
\textbf{MCC $\mu \pm \sigma$} &
\textbf{Interprétation}
\tabularnewline
\midrule

Aléatoire stratifiée appariée &
5 &
6 436 &
0,999260 $\pm$ 0,000353 &
0,998523 $\pm$ 0,000703 &
Performance très élevée lorsque des observations proches peuvent apparaître des deux côtés de la séparation entraînement--test.
\tabularnewline

Cinq holdouts par fichier ou scénario &
5 &
6 436 &
0,157827 $\pm$ 0,347377 &
0,146206 $\pm$ 0,301650 &
Forte dispersion principalement liée aux scénarios maintenus hors de l'entraînement.
\tabularnewline

\bottomrule
\end{tabularx}
\end{table}

Le contraste est important. Avec la séparation aléatoire stratifiée, le F1 moyen atteint 0,999260. Avec les holdouts par scénario, il tombe à 0,157827.

Comme les principaux paramètres sont contrôlés entre les deux variantes, le changement de stratégie de séparation constitue ici un élément déterminant de l'écart observé.

Il faut toutefois éviter une autre simplification. Les cinq valeurs du protocole strict ne correspondent pas à cinq répétitions indépendantes d'une même expérience où seul le hasard aurait changé. Chaque graine est associée à un fichier ou à un scénario différent maintenu hors de l'entraînement.

L'écart-type obtenu mesure donc en grande partie la différence de difficulté entre plusieurs problèmes de généralisation.

Cette expérience ne permet pas d'affirmer que tous les modèles échoueraient nécessairement sur CICIDS2017. Elle montre plus précisément que le Random Forest utilisé dans \logminer{} dépend fortement des caractéristiques présentes dans les fichiers et scénarios observés pendant son entraînement.

\begin{table}[H]
\centering
\caption{Analyse de l'instabilité du holdout CICIDS2017}
\label{tab:cicids-holdout-instability}
\scriptsize

\begin{tabularx}{\textwidth}{p{1cm}p{3.4cm}Xrrrr}
\toprule
\textbf{Graine} &
\textbf{Fichier maintenu hors de l'entraînement} &
\textbf{Attaques du fichier} &
\textbf{TP} &
\textbf{FN} &
\textbf{F1} &
\textbf{PR-AUC}
\tabularnewline
\midrule

42 &
Friday Afternoon DDoS &
DDoS = 61 194 &
2 553 &
1 447 &
0,779185 &
0,953532
\tabularnewline

43 &
Friday Afternoon PortScan &
PortScan = 12 488 &
20 &
3 980 &
0,009948 &
0,915857
\tabularnewline

44 &
Friday Morning &
Bot = 1 966 &
0 &
1 966 &
0,000000 &
0,315521
\tabularnewline

45 &
Thursday Afternoon Infiltration &
Infiltration = 32 &
0 &
32 &
0,000000 &
0,008821
\tabularnewline

46 &
Thursday Morning WebAttacks &
Brute Force = 1 507 ; XSS = 652 ; SQL Injection = 21 &
0 &
2 180 &
0,000000 &
0,339342
\tabularnewline

\bottomrule
\end{tabularx}
\end{table}

\noindent\footnotesize
\textit{Note : les effectifs de la colonne ``attaques du fichier'' décrivent le fichier source complet ; les TP et FN sont calculés sur l'échantillon plafonné utilisé pour le benchmark.}
\normalsize

Le détail par scénario permet de mieux comprendre l'écart-type particulièrement élevé du F1.

Le modèle conserve une capacité de généralisation partielle lorsque le scénario DDoS est maintenu hors de l'entraînement, avec un F1 de 0,779185. Sur PortScan, la performance devient très faible, avec un F1 de 0,009948. Aucun vrai positif n'est observé dans les échantillons testés pour Bot, Infiltration et WebAttacks.

L'instabilité ne s'explique donc pas principalement par le choix numérique des graines. Elle provient surtout du fait que les graines correspondent ici à des scénarios d'attaque différents, dont certains sont beaucoup plus éloignés des données observées lors de l'entraînement.

\begin{figure}[H]
\centering
\includegraphics[width=\textwidth]{fig_supervised_models_f1.png}
\caption{Scores F1 des modèles supervisés.}
\label{fig:supervised}
\end{figure}

\section{Évaluation non supervisée}

Dans les expériences non supervisées, les sorties des modèles sont systématiquement considérées comme des anomalies candidates. Cette règle concerne notamment les données Windows et les exports Wazuh utilisés dans les résultats principaux.

Une anomalie candidate indique qu'un événement ou un groupe d'événements présente un comportement inhabituel selon les caractéristiques et le modèle utilisés. Elle ne constitue pas, à elle seule, la preuve d'une attaque. Une interprétation de sécurité nécessite un contexte supplémentaire, une corrélation avec d'autres signaux ou une validation par un analyste.

Les expériences sur HDFS et BGL demandent une précaution supplémentaire.

Lorsqu'un modèle non supervisé est ajusté et évalué sur le même ensemble, ou lorsque son taux de contamination est directement dérivé du taux réel d'anomalies du jeu de données, le résultat mesure surtout la séparabilité interne de l'échantillon. Il ne représente pas encore une capacité de généralisation indépendante.

Une évaluation plus rigoureuse suppose de séparer les données en ensembles distincts. Les blocs HDFS ou les périodes BGL doivent être répartis entre entraînement, validation et test. Le modèle doit être ajusté uniquement à partir de l'entraînement, tandis que le seuil doit être déterminé sans utiliser les labels de l'ensemble de test. L'évaluation finale peut alors être effectuée sur des données jamais utilisées pendant le réglage.

\begin{figure}[H]
\centering
\includegraphics[width=0.92\textwidth]{fig_validation_selection_f1.png}
\caption{Comparaison des meilleurs détecteurs par jeu de données de validation.}
\label{fig:validation}
\end{figure}

\section{Comparaison avec Wazuh}

Les exports Wazuh fournissent un cas intéressant, car les événements ont déjà été enrichis par un outil HIDS/SIEM opérationnel.

L'objectif de l'expérience n'est pas de démontrer que \logminer{} remplace Wazuh. Le prototype reçoit plutôt les événements produits par la plateforme et applique une couche analytique supplémentaire afin d'identifier des comportements rares parmi ces signaux.

Sur l'ensemble traité, 122 563 événements Wazuh sont analysés et 3 676 sont signalés comme anomalies candidates par le mécanisme non supervisé.

Le recouvrement entre les deux sources de signalement permet ensuite d'observer dans quelle mesure certaines anomalies proposées par \logminer{} correspondent à des événements déjà considérés comme pertinents par les mécanismes de sécurité de Wazuh.

\begin{figure}[H]
\centering
\includegraphics[width=0.92\textwidth]{fig_wazuh_logminer_overlap.png}
\caption{Recouvrement entre les signaux Wazuh et les anomalies candidates de \logminer{}.}
\label{fig:wazuh}
\end{figure}

Cette comparaison doit rester qualitative lorsque les labels de vérité terrain ne sont pas disponibles. Une absence de recouvrement ne signifie pas nécessairement que \logminer{} produit un faux positif, pas plus qu'un recouvrement ne suffit à confirmer une attaque.

\section{Faux positifs}

Les faux positifs constituent une difficulté importante pour tout système de détection. Une méthode trop sensible peut produire de nombreux signaux et augmenter inutilement la charge de travail de l'analyste.

Lorsque les jeux de données disposent de labels fiables et que les métriques nécessaires sont disponibles, le prototype calcule des indicateurs permettant d'estimer cette erreur.

La situation est différente pour certaines évaluations non supervisées. En l'absence de vérité terrain suffisamment fiable, un événement signalé comme anomalie ne peut pas être classé automatiquement comme vrai ou faux positif.

Les taux présentés doivent donc être interprétés uniquement pour les expériences où la comparaison avec des labels ou une référence appropriée est possible.

\begin{figure}[H]
\centering
\includegraphics[width=0.92\textwidth]{fig_false_positive_rates.png}
\caption{Faux positifs pour 1 000 lignes selon les métriques disponibles.}
\label{fig:false-positive}
\end{figure}

\section{Benchmark local par micro-lots}

Une campagne spécifique a été réalisée afin de mesurer le temps nécessaire à l'exécution du workflow local.

Le benchmark utilise l'endpoint FastAPI \texttt{/run/discovered}. Dix cycles sont exécutés, avec 8 537 lignes traitées à chaque cycle.

La durée moyenne du workflow est de 8,2012 secondes. La durée minimale observée est de 3,1672 secondes, tandis que la valeur maximale atteint 15,3289 secondes.

\begin{table}[H]
\centering
\caption{Benchmark local par micro-lots}
\label{tab:realtime}

\begin{tabular}{lr}
\toprule
\textbf{Indicateur} &
\textbf{Valeur}
\tabularnewline
\midrule

Cycles réussis &
10
\tabularnewline

Lignes par cycle &
8 537
\tabularnewline

Durée moyenne du workflow &
8,2012 s
\tabularnewline

Durée minimale du workflow &
3,1672 s
\tabularnewline

Durée maximale du workflow &
15,3289 s
\tabularnewline

\bottomrule
\end{tabular}
\end{table}

Ces valeurs décrivent un traitement périodique de micro-lots dans l'environnement utilisé pour l'expérience. Elles ne constituent pas une démonstration de traitement en temps réel strict.

En particulier, l'expérience ne mesure pas la capacité du système à absorber continuellement un flux de très haut débit avec une contrainte de latence garantie. Elle permet plutôt d'obtenir un ordre de grandeur du temps nécessaire au workflow local sur une charge contrôlée.

\begin{figure}[H]
\centering
\includegraphics[width=0.92\textwidth]{fig_realtime_workflow_latency.png}
\caption{Durée du workflow local par micro-lots sur 10 cycles.}
\label{fig:latency}
\end{figure}

\section{Campagne CPU/RAM}

Une campagne de 30 cycles complète les mesures de latence en observant la consommation de ressources.

Elle mesure la durée du workflow, l'utilisation du CPU ainsi que la mémoire RAM consommée par l'API/orchestrateur et par le processus principal de \logminer{}.

\begin{table}[H]
\centering
\caption{Campagne de mesure des ressources sur plusieurs cycles}
\label{tab:resources}
\scriptsize

\begin{tabular}{lrrrr}
\toprule

\textbf{Agent} &
\textbf{Cycles} &
\textbf{Workflow moy.} &
\textbf{CPU machine moy.} &
\textbf{RAM moy.}
\tabularnewline

\midrule

API / orchestrateur &
30 &
9,3300 s &
7,45\% &
187,82 MB
\tabularnewline

Processus \logminer{} &
30 &
9,3300 s &
0,41\% &
495,41 MB
\tabularnewline

\bottomrule
\end{tabular}
\end{table}

Sur les 30 cycles, la durée moyenne du workflow est de 9,3300 secondes.

L'API/orchestrateur utilise en moyenne 7,45\% du CPU de la machine et 187,82 MB de RAM. Le processus \logminer{} présente une utilisation CPU moyenne de 0,41\% dans la mesure considérée, pour une consommation moyenne de mémoire de 495,41 MB.

Ces mesures décrivent le comportement du prototype dans la configuration matérielle et logicielle utilisée pour la campagne. Elles ne doivent donc pas être interprétées comme des valeurs universelles indépendantes de la machine, du système d'exploitation ou de la charge traitée.

\begin{figure}[H]
\centering
\includegraphics[width=0.92\textwidth]{fig_resource_campaign_multicycle.png}
\caption{Évolution de l'utilisation du CPU et de la RAM pendant la campagne multicycle.}
\label{fig:resources}
\end{figure}

\section{Campagne d'exécution parallèle}

Une campagne complémentaire a été menée pour observer le comportement de l'exécution parallèle.

L'expérience utilise trois workers et un fichier CSV synthétique labellisé contenant 500 événements à chaque cycle. Les cinq cycles prévus se terminent correctement.

La durée moyenne du workflow est de 3,6232 secondes et la durée maximale de 6,0172 secondes. Le pic moyen d'utilisation du CPU de la machine atteint 17,1725\%, tandis que la consommation maximale moyenne de RAM est de 176,89 MB.

\begin{table}[H]
\centering
\caption{Campagne de mesure des ressources en mode parallèle}
\label{tab:parallel-resources}
\scriptsize

\begin{tabular}{lr}
\toprule
\textbf{Indicateur} &
\textbf{Valeur}
\tabularnewline
\midrule

Cycles réussis &
5/5
\tabularnewline

Événements par cycle &
500
\tabularnewline

Workers parallèles &
3
\tabularnewline

Durée moyenne du workflow &
3,6232 s
\tabularnewline

Durée maximale du workflow &
6,0172 s
\tabularnewline

CPU machine maximal moyen &
17,1725\%
\tabularnewline

RAM maximale moyenne &
176,89 MB
\tabularnewline

\bottomrule
\end{tabular}
\end{table}

Cette campagne confirme surtout que le mode parallèle peut être exécuté correctement avec plusieurs workers dans l'environnement local.

Elle ne doit pas être comparée directement au benchmark précédent comme s'il s'agissait d'une accélération mesurée à charge identique, puisque les volumes, les données et les conditions expérimentales ne sont pas les mêmes.

\begin{figure}[H]
\centering
\includegraphics[width=0.92\textwidth]{fig_parallel_resource_campaign.png}
\caption{Pics d'utilisation du CPU observés pendant la campagne parallèle locale.}
\label{fig:parallel-resources}
\end{figure}

\section{Expérimentations HDFS et BGL}

Les résultats obtenus sur HDFS et BGL sont volontairement séparés des résultats supervisés principaux.

Ces deux jeux de données possèdent une composante séquentielle importante. Une anomalie peut dépendre d'une succession de messages, d'un bloc, d'une fenêtre temporelle ou d'un changement dans la fréquence d'un modèle de message. Une analyse purement ligne par ligne risque donc de perdre une partie de cette information.

Deux niveaux d'expérimentation sont présentés.

Le premier correspond aux essais exploratoires effectués à partir de caractéristiques légères extraites ligne par ligne. Ces expériences sont utiles pour vérifier l'intégration des données dans le pipeline, mais elles ne constituent pas une validation forte de la capacité de généralisation.

\begin{table}[H]
\centering
\caption{Résultats exploratoires sur HDFS et BGL}
\label{tab:hdfs-bgl-exploratoires}
\scriptsize

\begin{tabularx}{\textwidth}{p{2.5cm}p{3.4cm}p{3.3cm}X}
\toprule

\textbf{Famille} &
\textbf{Résultat observé} &
\textbf{Interprétation retenue} &
\textbf{Limite principale}
\tabularnewline

\midrule

HDFS &
F1 autour de 0,599333--0,600333 &
Séparabilité partielle avec des caractéristiques légères extraites ligne par ligne. &
Les anomalies HDFS dépendent souvent de séquences ou de modèles de messages ; ce protocole ne démontre pas une généralisation indépendante.
\tabularnewline

BGL &
F1 autour de 0,994333 &
Forte séparabilité dans l'échantillon traité et validation de l'intégration au pipeline. &
Le score est sensible au protocole \textit{in-sample} et au choix du seuil ; il ne doit pas être interprété comme une validation robuste hors échantillon.
\tabularnewline

\bottomrule
\end{tabularx}
\end{table}

Ces premiers résultats ont donc principalement une valeur exploratoire. Ils montrent que \logminer{} peut importer et traiter ces journaux, mais ils ne permettent pas à eux seuls de conclure à la robustesse des détecteurs sur des données indépendantes.

Une seconde expérimentation a été mise en place afin de renforcer le protocole.

Les données sont d'abord séparées entre entraînement et test. Pour HDFS, la partition respecte les identifiants de blocs afin d'éviter qu'un même bloc apparaisse dans les deux ensembles. Pour BGL, la séparation suit une logique chronologique stratifiée.

Les messages sont ensuite enrichis à l'aide de Drain3. Le pipeline exploite les modèles de messages, leur fréquence, leur rareté, des fenêtres temporelles ainsi que plusieurs informations liées au contexte de la source.

Les modèles sont ajustés sur l'ensemble d'entraînement puis évalués sur l'ensemble de test.

\begin{table}[H]
\centering
\caption{Validation Drain3 entraînement--test sur HDFS et BGL}
\label{tab:hdfs-bgl-drain3-train-test}
\scriptsize

\begin{tabularx}{\textwidth}{p{2.6cm}p{3.4cm}rrrrX}
\toprule

\textbf{Famille} &
\textbf{Modèle} &
\textbf{Précision} &
\textbf{Rappel} &
\textbf{F1} &
\textbf{Exact.} &
\textbf{Interprétation méthodologique}
\tabularnewline

\midrule

HDFS &
Histogramme entraînement--test &
0,653333 &
0,652246 &
0,652789 &
0,652789 &
Gain modéré par rapport aux caractéristiques ligne par ligne ; la famille reste difficile à traiter.
\tabularnewline

HDFS &
Isolation Forest entraînement--test &
0,638333 &
0,637271 &
0,637802 &
0,637802 &
Résultat proche du meilleur histogramme, mais avec encore un nombre important de faux positifs et de faux négatifs.
\tabularnewline

BGL &
Isolation Forest entraînement--test &
1,000000 &
1,000000 &
1,000000 &
1,000000 &
Séparation parfaite observée sur cette partition locale ; ce résultat demande des répétitions supplémentaires avant toute généralisation.
\tabularnewline

BGL &
Z-score / IQR / histogramme / autoencodeur &
0,973333 &
0,973333 &
0,973333 &
0,973333 &
Forte séparabilité confirmée, avec quelques erreurs résiduelles selon le détecteur.
\tabularnewline

\bottomrule
\end{tabularx}
\end{table}

La seconde expérience améliore la cohérence méthodologique de l'évaluation, notamment en rapprochant la représentation des données de celles utilisées dans les travaux consacrés aux journaux séquentiels.

L'amélioration reste modérée sur HDFS. Le meilleur F1 présenté atteint 0,652789 avec l'approche par histogramme, contre des valeurs proches de 0,60 dans les essais ligne par ligne.

Sur BGL, Isolation Forest obtient un score parfait sur la partition locale utilisée. Ce résultat indique une très forte séparabilité dans cette configuration, mais il doit être considéré avec prudence. Une seule séparation favorable ne suffit pas à démontrer une performance généralisable.

Les résultats obtenus directement sur des fichiers de test sans séparation stricte entre entraînement et évaluation sont donc conservés uniquement comme expériences secondaires. Les résultats Drain3 avec séparation entraînement--test constituent la référence principale pour cette partie du mémoire.

\section{Campagne Redis des agents intelligents}

Une campagne d'endurance d'une durée de six heures a été menée afin d'évaluer le fonctionnement des agents intelligents multitâches avec Redis Streams.

Le scénario utilise trois workers principaux ainsi qu'un worker chargé de la reprise. À chaque itération, une panne est volontairement simulée avant l'acquittement d'une tâche afin de vérifier la capacité du système à récupérer les messages qui restent en attente.

Une campagne plus courte, identifiée par \texttt{redis-campaign-20260717161257}, avait auparavant validé 150 tâches. Elle est considérée comme une prévalidation. La campagne de six heures constitue l'expérience principale.

\begin{table}[H]
\centering
\caption{Campagne Redis d'endurance des agents intelligents}
\label{tab:intelligent-redis-long}
\scriptsize

\begin{tabular}{lr}
\toprule
\textbf{Indicateur} &
\textbf{Valeur}
\tabularnewline
\midrule

Itérations terminées & 709 \tabularnewline
Tâches mises en file & 8 508 \tabularnewline
Tâches uniques terminées & 8 508 \tabularnewline
Échecs & 0 \tabularnewline
Pannes simulées avant acquittement & 709 \tabularnewline
Tâches finales en attente & 0 \tabularnewline
Perte estimée & 0 \tabularnewline
Durée observée & 21 613,8566 s \tabularnewline
Débit observé & 0,3936 tâches/s \tabularnewline
Latence p95 & 5,9879 s \tabularnewline
Latence p99 & 8,7987 s \tabularnewline

\bottomrule
\end{tabular}
\end{table}

Au total, 8 508 tâches uniques ont été mises en file et les 8 508 ont été terminées. Aucune tâche finale ne reste en attente et aucune perte n'est observée dans les conditions de l'expérience.

Les 709 pannes simulées donnent lieu à 709 reprises par le worker prévu à cet effet.

La répartition des tâches confirme également que les différents workers participent effectivement à l'exécution. \texttt{redis-agent-1} termine 2 654 tâches, \texttt{redis-agent-2} en exécute 2 602 et \texttt{redis-agent-3} en traite 2 543. Le \texttt{redis-recovery-agent} prend en charge les 709 tâches abandonnées lors des pannes simulées.

La charge est également équilibrée entre les trois types de tâches : 2 836 exécutions de \texttt{parse.logs}, 2 836 de \texttt{route.model} et 2 836 de \texttt{discover.logs}.

Cette campagne fournit donc une preuve expérimentale du fonctionnement distribué à l'échelle locale multiprocessus. Elle montre notamment que plusieurs workers peuvent partager le même flux de tâches et qu'un message non acquitté peut être repris.

Elle ne constitue pas pour autant une preuve de déploiement distribué industriel. Les expériences restent réalisées dans un environnement local et ne traitent pas l'ensemble des problèmes d'un système multi-machine réel, notamment la sécurité du bus, les partitions réseau, la réplication ou la tolérance à la panne de plusieurs nœuds.

Une comparaison contrôlée a également été réalisée entre une exécution monolithique, une exécution à base d'agents et une variante dans laquelle une panne est simulée.

\begin{table}[H]
\centering
\caption{Comparaison contrôlée entre architecture monolithique et architecture à base d'agents}
\label{tab:controlled-monolith-agents}
\scriptsize

\begin{tabularx}{\textwidth}{p{3cm}rrrrrX}
\toprule

\textbf{Mode} &
\textbf{Tâches} &
\textbf{Durée} &
\textbf{Débit} &
\textbf{CPU moy.} &
\textbf{RAM max.} &
\textbf{Interprétation}
\tabularnewline

\midrule

Monolithique &
60/60 &
36,6815 s &
1,6357 t/s &
121,339 &
189,074 MB &
Référence séquentielle contrôlée.
\tabularnewline

Agents &
60/60 &
38,1803 s &
1,5715 t/s &
165,746 &
234,777 MB &
Surcoût d'orchestration visible sur un faible volume de tâches.
\tabularnewline

Agents avec reprise &
60/60 &
39,3000 s &
1,5267 t/s &
117,598 &
189,395 MB &
Une panne simulée, une tâche reprise et aucune tâche finale en attente.
\tabularnewline

\bottomrule
\end{tabularx}
\end{table}

Les résultats ne montrent pas d'amélioration du débit liée à l'utilisation des agents sur ce faible volume.

Le mode monolithique traite 1,6357 tâche par seconde, contre 1,5715 pour le mode utilisant les agents. Lorsque la reprise après panne est activée, le débit descend à 1,5267 tâche par seconde.

L'architecture multiagent introduit donc ici un coût d'orchestration mesurable.

Cette absence de gain de débit ne remet pas en cause la fonction architecturale du mécanisme. La variante avec reprise montre qu'une tâche abandonnée peut être récupérée sans rester en attente à la fin du scénario.

Les bénéfices revendiqués concernent donc surtout la séparation des responsabilités, l'observabilité, la possibilité de distribuer localement les traitements et la reprise après panne. La campagne ne permet pas de présenter les agents comme une optimisation systématique des performances.

\section{Validation progressive de l'architecture d'exécution}

Les différents scénarios d'exécution permettent également de suivre l'évolution du prototype depuis une chaîne centralisée vers une organisation comportant plusieurs agents.

Cette progression ne constitue pas une étude d'ablation scientifique stricte. Les différentes configurations ne manipulent pas exactement les mêmes volumes ni les mêmes charges de travail.

Il serait donc incorrect de comparer directement leurs débits afin de conclure à la supériorité d'une architecture.

L'objectif est plus limité : vérifier successivement que le traitement peut fonctionner dans un pipeline centralisé, qu'un agent peut prendre en charge plusieurs types de tâches, que plusieurs agents peuvent partager une file et qu'une exécution Redis peut récupérer une tâche après une panne simulée.

\begin{table}[H]
\centering
\caption{Validation progressive des agents intelligents}
\label{tab:intelligent-agents-ablation}
\scriptsize

\begin{tabularx}{\textwidth}{p{3.2cm}rrrrX}
\toprule

\textbf{Scénario} &
\textbf{Tâches} &
\textbf{Échecs} &
\textbf{Durée} &
\textbf{Débit} &
\textbf{Conclusion}
\tabularnewline

\midrule

Pipeline centralisé &
2 &
0 &
n/a &
n/a &
Référence non distribuée.
\tabularnewline

Agent unique multitâche &
3 &
0 &
n/a &
n/a &
Un agent exécute la découverte, le parsing et le routage.
\tabularnewline

Agents locaux multiples &
9 &
0 &
6,0688 s &
1,483 &
Partage d'une file locale entre plusieurs agents.
\tabularnewline

Agents Redis -- prévalidation &
150 &
0 &
102,4202 s &
1,4646 &
Distribution locale multiprocessus validée.
\tabularnewline

Agents Redis -- endurance 6 h &
8 508 &
0 &
21 613,8566 s &
0,3936 &
Distribution locale multiprocessus et reprise répétée validées.
\tabularnewline

\bottomrule
\end{tabularx}
\end{table}

La progression montre que les différents mécanismes prévus par l'architecture peuvent être activés successivement.

Le résultat le plus important n'est donc pas le débit brut de chaque ligne du tableau. Il réside dans le passage d'un traitement centralisé à une file partagée, puis à un bus Redis avec plusieurs processus et récupération répétée de tâches non acquittées.

\section{Robustesse multiformat}

La robustesse du parsing est testée sur plusieurs formats de journaux ainsi que sur des entrées incomplètes ou inhabituelles.

L'objectif principal est de vérifier que le pipeline ne supprime pas silencieusement une entrée simplement parce qu'elle ne correspond pas parfaitement à l'un des formats connus.

Lorsqu'une ligne peut être interprétée, les champs disponibles sont extraits et normalisés. Lorsqu'elle reste partiellement inconnue ou corrompue, le système conserve le message original ainsi qu'un statut permettant de retrouver le problème.

Cette approche ne garantit pas que toute entrée inconnue pourra être analysée correctement. Elle vise plutôt à éviter une perte silencieuse d'information, ce qui est particulièrement important dans une chaîne de sécurité.

\begin{figure}[H]
\centering
\includegraphics[width=0.92\textwidth]{fig_robustness_multiformat.png}
\caption{Robustesse du parsing multiformat.}
\label{fig:robustness}
\end{figure}

\section{Synthèse du chapitre}

Les expériences présentées dans ce chapitre permettent d'évaluer plusieurs dimensions du prototype plutôt que de réduire son appréciation à un score unique.

Sur les données supervisées, les résultats exploratoires sont complétés par des partitions plus strictes. Le cas de CICIDS2017 met clairement en évidence la sensibilité du modèle au protocole de séparation : le F1 passe de 0,999260 dans la comparaison aléatoire appariée à une moyenne de 0,157827 lorsque des scénarios ou fichiers entiers sont maintenus hors de l'entraînement.

Linux/auth présente une baisse plus modérée avec un F1 strict de 0,826251, tandis que CIC-DDoS2019 conserve un F1 élevé de 0,997577 dans le protocole de holdout retenu.

Les expériences HDFS et BGL distinguent désormais les essais exploratoires des évaluations utilisant Drain3, des fenêtres temporelles et une séparation entraînement--test. Cette distinction permet de mieux tenir compte de la nature séquentielle de ces journaux sans attribuer aux scores exploratoires une portée qu'ils ne possèdent pas.

Les mesures système apportent une autre dimension à l'évaluation. Elles documentent le temps d'exécution du workflow, la consommation CPU/RAM et le fonctionnement d'une exécution parallèle locale.

Enfin, la campagne Redis de six heures constitue la principale validation de la dimension distribuée à l'échelle locale. Elle traite 8 508 tâches uniques, simule 709 pannes avant acquittement et termine sans tâche finale en attente. La comparaison avec le mode monolithique montre cependant que cette architecture introduit un surcoût sur de faibles volumes de tâches. Son intérêt se situe donc davantage dans l'organisation du traitement et la capacité de reprise que dans un gain de débit démontré.

Ces résultats constituent la base expérimentale de la discussion. Leur portée, leurs limites ainsi que les conséquences méthodologiques des différents protocoles sont examinées dans les chapitres suivants.
